\section{The DP4 library: what we are shipping}\label{sec:dp4}
This section inventories the components we contribute to the DP4 assay run and what each one is
for. DP4 as a whole is the larger fourth PepSeq library; the pieces assembled here are the
designed and control constructs this project adds to it. Each component targets one of the three
questions from Section~\ref{sec:intro}, and all of them are emitted in the same constant-length
construct (Section~\ref{sec:tiling}) and the same 8-column annotated format, so they can be merged into
one ordered synthesis file (Section~\ref{sec:assembly}).

\begin{table}[h]\centering\small
\begin{tabular}{llcl}
\toprule
component & what it is & constructs & status \\
\midrule
C1 known antibody, whole epitope & full-epitope scaffolds, top-\num{20}/epitope & \num{1120} & ranked \\
C2 known antibody, single island & 87 island contigs, each island alone & \num{1600} & ranked, per island \\
C3 scaffolds, polyclonal tiling & top-\num{10}/window, no antibody & \num{4390} & selected, ungated \\
C4 linear tiled-30mer controls & bare tiled peptides & \num{2033} & built \\
C5 metric-space sampling & designs spread across metrics & \num{2715} & sampled \\
C6 scaffolded-epitope controls & island$\to$Ala $+$ scaffold-disruption & \num{2325} & built \\
8VDL PfEMP1 conserved epitope & two epitope definitions, top-\num{10} each & \num{20} & ranked \\
\bottomrule
\end{tabular}
\caption{The DP4 components assembled here, at the shipped depth. C1 and C2 are the two halves of the
islands experiment (whole epitope vs.\ one island at a time); C6 tests the same question by mutation.
C1--C3 are selected (ranked on the composite, Section~\ref{sec:selection}); the shipped depth is
top-\num{20} per group for C1/C2 and top-\num{10} per window for C3. C5 is sampled at \num{3000} and C4
built at \num{2034} tiles (Section~\ref{sec:tiling}); the counts in the table are the shipped totals after
the 2026-07-21 cull (Section~\ref{sec:assembly}). Together with the 8VDL arm this is \num{14203}
shipped constructs. The delivered \texttt{dp4\_library.csv} additionally carries the \num{21759} filter-passing
LX PfEMP1/EPCR minibinders --- a separate de-novo binder arm, not scaffolded or scored here --- for a
single-file view of the whole library (\num{35962} rows); the \num{14203} above is the episcaf portion
the scoring concerns, while the oligo encoding and order file cover the whole \num{35962}.}
\label{tab:dp4components}
\end{table}

\subsection{C1 --- known-antibody epitopes, whole}\label{sec:dp4c1}
For each mAb epitope we ship the best few scaffolds that present the \emph{complete} epitope ---
all of its islands held together in their native geometry, the way the antibody actually sees it.
These are the designs whose binding we can score against a ground truth, since the antibody is
known; they continue the DP3 set that gave us the binding data of Section~\ref{sec:experimental}
and are the comparator for the single-island arm below. Selection ranks on the re-weighted
composite with no hard gate (Section~\ref{sec:selection}); because the antibody is known, the
accessibility term is the \emph{real} clash \texttt{af3\_n\_clash\_res} rather than the cylinder
surrogate. Whole-epitope presentation is straightforward for the \num{13} single-island epitopes
but hard for multi-island ones --- preserving the relative geometry of separated islands is
exactly where scaffolding struggles (Section~\ref{sec:intro}, the \texttt{2XQB} failure) --- which
is the motivation for shipping C2 alongside it. The C1 design set --- how it is built at the native
\num{103}-mer length, its yield, and its selection --- is characterized in Section~\ref{sec:c1set}; C5
and C6 derive from that same pool.

\subsection{C2 --- known-antibody epitopes, one island at a time}\label{sec:dp4c2}
The single-island component scaffolds each island of the \num{46} multi-island epitopes
\emph{individually}: one island per design, presented alone. This is the direct test of Q1
(Section~\ref{sec:islands}) --- does an antibody need both islands of its epitope, or does one
suffice? Shipping C1 (whole epitope) and C2 (each island alone) for the same antibodies lets the
assay deconvolve which island, if any, carries the binding. The per-island production run
(Section~\ref{sec:perisland}) covers \num{87} scaffolded islands (\num{5} size-1 islands are
skipped as un-presentable); at the shipped top-\num{20} per island (after the 56-mAb exclusion) that
is \num{1600} shipped constructs (Table~\ref{tab:dualisland}). Selection is identical to C1 (composite, no
gate, real-clash accessibility), ranked per (\texttt{id}, island) so each island competes on its own.

\subsection{C3 --- scaffolded epitopes from polyclonal tiling}\label{sec:dp4c3}
For antigens with no known antibody we tile the whole protein into overlapping windows and
scaffold each one (Section~\ref{sec:tiling}), then keep the top \num{10} designs per window by the
composite. This is the polyclonal arm (Q3): we do not know where the real epitope is, so we cover
the antigen densely and let the assay find it. This is the one component where the antibody-clash
filter cannot be computed, so accessibility is judged entirely by the native-aware cylinder
surrogate (Section~\ref{sec:noab}) --- the metric the DP3 binding data singled out as the
strongest predictor (Section~\ref{sec:whatpredicts}), which is why it now carries the most weight
in selecting these designs. Across the three tiled antigens this is \num{1910} (1D2K) $+$ \num{1560}
(4WAT) $+$ \num{920} (6M0J) $=$ \num{4390} designs: ten per scorable window, over the \num{191}, \num{156},
and \num{92} windows that clear the crystal-coverage check (the ``epitopes kept'' column of
Table~\ref{tab:tilecount}, which itself lists the top-\num{5} counts).

\paragraph{Why C3 is shipped deep.} The C3 windows are \num{12}-mers stepping by \num{2} residues, so
neighbouring tiles are highly redundant --- adjacent windows share \num{10} of \num{12} residues, which
argues for a shallow per-window cut. We nonetheless ship top-\num{10} rather than a shallow depth,
because that redundancy does not protect against the failure mode we actually face here: the polyclonal
designs have the weakest accessibility (clash) distribution of any component, so the per-design success
probability is low, and overlapping tiles can fail together. This is also the arm with the most to gain
--- even a handful of hits in the no-antibody setting would be the first demonstration that the
approach works without a known antibody --- so the spare capacity in the peptide budget is spent here,
maximizing the number of independent chances per window.

\subsection{C4 --- linear tiled-30mer controls}\label{sec:dp4c4}
The linear controls are the bare epitope peptides the scaffolded designs are meant to beat:
unscaffolded 30-mers (model \texttt{(none)}) tiled across the antigen, presenting the sequence
with none of the conformational constraint a scaffold imposes. They are the baseline for the
central scaffold-versus-linear question (Section~\ref{sec:intro}): if conformational presentation
matters, a scaffolded design should bind its antibody more reliably than the corresponding linear
tile. We tile the gap-free PDB FASTA chain at a 30-residue window, step 6
(Section~\ref{sec:tiling}), giving \num{2034} tiles over \num{59} antigens --- the \num{56}
known-antibody targets plus the \num{3} polyclonal antigens
(\texttt{data/libraries/dp4\_tiled30mers\_fasta.csv}). The three non-standard cases excluded from the
antibody arms (\texttt{2h32} the pre-B-cell receptor, \texttt{4xwo} low-yield, \texttt{7a3t} too-small;
Section~\ref{sec:dp4c6}) are dropped here as well, so the linear controls stay consistent with the
scaffold arms they are meant to be compared against. This component is built.

\subsection{C5 --- metric-space sampling}\label{sec:dp4c5}
The DP3 binding data is a weak prior precisely because every assayed design already passed the
four filters, so the metrics barely varied (Section~\ref{sec:experimental}). C5 fixes that by
design: instead of shipping only the top designs, we deliberately sample designs \emph{spread
across} the metric space --- including ones the current filters would reject --- so that epitope
RMSD, accessibility, and PAE each take a wide range of values. Accessibility is sampled two ways, split evenly: half the designs span the \emph{real} AlphaFold3
clash (available because the antibody is known) and half the native-aware cylinder surrogate, so the
titration calibrates both the ground-truth accessibility we have for the antibody set and the surrogate
we must rely on where no antibody is known. Reading binding off this spread is what lets us learn which
metrics actually predict enrichment and finally fit the composite's weights (Q2,
Section~\ref{sec:open}), rather than re-deriving the filters we started from. The sampling is implemented
(\texttt{scripts/stage06\_sample\_c5.py}): a farthest-point (max--min) spread over four standardized
scoring axes --- epitope RMSD, epitope PAE, overall RMSD, and one accessibility axis --- sampled per mAb
over the \num{56} assayable epitopes, with the accessibility axis being the real clash for one
(disjoint) half of the \num{3000} designs and the cylinder for the other (about \num{54} per mAb).

The sample reaches deep into the tails the filters reject, covering most of each axis's full range
(Table~\ref{tab:c5span}, Fig.~\ref{fig:c5span}).

\begin{table}[h]\centering\small
\begin{tabular}{lccc}
\toprule
axis & pool range & sample range & span \\
\midrule
epitope RMSD (\si{\angstrom}) & 0.16--60.4 & 0.20--39.6 & \SI{65}{\percent} \\
epitope PAE                   & 1.0--21.0  & 1.1--21.0  & \SI{100}{\percent} \\
overall RMSD (\si{\angstrom}) & 0.31--45.5 & 0.35--44.5 & \SI{98}{\percent} \\
AF3 clash (residues)          & 0--85      & 0--81      & \SI{95}{\percent} \\
native-aware cylinder         & 0--57      & 0--51      & \SI{93}{\percent} \\
\bottomrule
\end{tabular}
\caption{C5 titration coverage --- the fraction of each axis's full pool range the sample spans (span
$=$ sample range $/$ pool range). All axes are near-fully spanned; epitope RMSD reads \SI{65}{\percent}
only because its pool maximum (\SI{60}{\angstrom}) is a handful of unfolded outliers the sample does not
chase, so the range ratio understates the coverage of the meaningful range.}
\label{tab:c5span}
\end{table}

\begin{figure}[h]\centering
\figorbox{c5_titration_coverage.png}{1.0}
\caption{C5 metric-space titration, one panel per sampled axis: the full whole-epitope pool (grey) vs the
\num{3000}-design farthest-point sample (blue). The sample tracks pool density on any single marginal ---
it spreads in the \emph{joint} space, not per-axis --- yet reaches the tails; the red line is a
stratified-uniform ``span this axis alone'' reference (flat by construction) showing what per-axis-uniform
coverage would look like. Each curve is density-normalized independently. Regenerated by
\texttt{episcaf\_analysis/viz/plot\_c5\_titration.py}.}
\label{fig:c5span}
\end{figure}

\subsection{C6 --- scaffolded-epitope controls}\label{sec:dp4c6}
C6 is the mutation-based read of the islands question (Q1), and it is the complement of C2: where C2
scaffolds each island alone to ask whether one island \emph{suffices}, C6 knocks out an island inside the
intact design to ask whether each island is \emph{necessary}. The unmutated designs are not part of this
component --- they already ship as C1 --- so C6 is the controlled comparison against them.

Mechanically, for the best \num{20} designs of each of the \num{56} mAbs we alanine-scan the epitope
islands: every residue of island~1 mutated to alanine, and (where the epitope has a second island) every
residue of island~2. Alanine strips a residue's side-chain chemistry while leaving the fold intact, so if
binding drops when island~1 is alanine-substituted the antibody requires island~1, and likewise for
island~2. The mutations are string substitutions on the existing designs' case-encoded sequence (epitope
residues uppercase, scaffold lowercase, each contiguous uppercase run an island), implemented in
\texttt{episcaf\_pipeline/scaffolded\_epitope\_controls/} (a port of the DP3 scheme). Over the best
designs of each of the \num{56} mAbs, with one island-mutant for single-island epitopes and two for
multi-island ones, plus the scaffold-disruption variant below, this is \num{3100} constructs at a top-20
base (\num{1980} island-alanine mutants $+$ \num{1120} scaffold-disruption controls, one per base design).
The shipped depth is top-\num{15} (\num{2325} constructs) --- C6 was the biggest block for what are
controls, so it was scaled down in the cull (Section~\ref{sec:assembly}).

\paragraph{Scaffold-disruption control.} A companion negative control mutates the \emph{scaffold}
(not the epitope) in each design: if binding is epitope-driven, perturbing the scaffold should leave
it intact, so a design that survives this but collapses under an island alanine-substitution is
behaving as intended. We insert a fixed structure-disrupting hexamer (\texttt{PPDDGG}) into four
lowercase (scaffold) windows, each required to lie at least four residues from any epitope position,
so the epitope is untouched and the perturbation is confined to the scaffold; placement is seeded for
reproducibility. We ship this four-window variant (\texttt{scaffoldMutX4}) --- the stronger
perturbation carried over from DP3 --- and drop the single-window one to save space. Because island
membership and the epitope-avoidance constraint are both read directly from the case-encoding, no
separate residue-index map is needed. When a scaffold cannot fit four disjoint windows, we
\emph{degrade the count} (\texttt{scaffoldMutX3}, \texttt{X2}, \texttt{X1}) rather than drop the control,
so a design loses its scaffold control only if it cannot place even a single hexamer $\ge\!4$ residues
from the epitope. In the shipped build the four-window arm covers \num{1034} designs, with \num{60} at
three windows, \num{25} at two, and \num{1} at a single window, so \emph{every} one of the \num{1120}
bases carries a scaffold control; the island-alanine arms likewise cover all bases.

\medskip
\noindent The controls follow the DP3 mutation scheme (reference and port in
\texttt{episcaf\_pipeline/scaffolded\_epitope\_controls/}); the DP4 changes are all-residue island
alanine (not every other) and the four-window scaffold variant only. The base set is the same top-$n$
per mAb used for C1.

\paragraph{The known-antibody set (\num{56}).} The antibody-based components (C1, C2, C5, C6) share one
exclusion set applied consistently: from the \num{59} DP3 mAb epitopes we drop three, leaving \num{56}.
Two are the assay drops of Section~\ref{sec:experimental} (\texttt{4xwo}, low antibody yield;
\texttt{7a3t}, a \num{4}-residue epitope too small to present). The third, \texttt{2h32}, is excluded
because it is not a conventional antibody:antigen complex at all --- it is the pre-B-cell receptor, which
engages its ligand differently, so it is not a valid test case for antibody-epitope scaffolding; it
entered the DP3 set through the inclusion criteria inadvertently. The polyclonal tiling (C3) is over
different antigens and is unaffected.

\subsection{8VDL --- a PfEMP1 conserved-epitope arm}\label{sec:dp4_8vdl}
One further arm scaffolds a conserved epitope from a different antigen entirely: the EPCR-binding
surface of the \emph{Plasmodium falciparum} PfEMP1 CIDR$\alpha$1.4 domain (crystal structure
\texttt{8VDL}; Reyes et al., \emph{Nature} 636:182--189, 2024). PfEMP1 is the hypervariable parasite
surface antigen that drives severe malaria and escapes immunity by antigenic variation, but the
residues its CIDR domain uses to grip the host receptor EPCR cannot vary freely --- so a construct that
presents that conserved contact surface is a candidate for eliciting \emph{broadly} reactive,
variant-transcending antibodies. The paper identifies the conserved, EPCR-binding-critical residues on
which its broadly reactive antibodies depend: the double-phenylalanine motif \texttt{F655}/\texttt{F656}
and a supporting glutamate \texttt{E666}. We scaffold two epitope definitions below --- the contiguous
chain~C window covering the antibody's contact surface, and those conserved hotspots alone. Because the
crystal also contains the cognate C7 antibody (chains H/L), this is a known-antibody target and the real
clash term applies, exactly as for C1/C2.

We scaffold it two ways and keep the best \num{10} of each: \texttt{epitope} fixes the whole
contiguous contact window \texttt{C652--C673} (\num{22} residues, spanning all \num{13} residues with a
heavy atom within \SI{4}{\angstrom} of the Fab --- the strongest constraint, the direct analog of C1),
while \texttt{hotspots} fixes only the three functional residues \texttt{F655}/\texttt{F656}/\texttt{E666}
and lets the design build everything else around them (a minimal ``hotspot graft''). Both are scaffolded
into the same \num{103}-mer construct through the same
RFdiffusion3~$\rightarrow$~ProteinMPNN~$\rightarrow$~AlphaFold3 pipeline and ranked under the same
soft-gate composite the antibody arms use (Section~\ref{sec:composite}; the custom consolidation
\texttt{dp4\_8vdl/scripts/07\_consolidate.py} aligns each predicted epitope onto the native chain~C
frame and scores the real H/L clash there, so accessibility is the true Fab clash, not the cylinder
surrogate). The fixed atoms carry the residues' native crystal coordinates, so the hotspot arm preserves
their three-dimensional arrangement however the construct spaces them. The two arms give a clean,
testable contrast --- and, as it turns out, one of them fails informatively. The whole-epitope designs
are antibody-\emph{accessible}: their top \num{10} recover the epitope well (epitope RMSD
\SIrange{0.97}{1.36}{\angstrom}) while leaving the Fab surface largely unobstructed (\numrange{1}{2}
clashing residues). The minimal hotspot grafts recover the three residues almost exactly (epitope RMSD
\SIrange{0.04}{0.11}{\angstrom}) but bury them under scaffold that would block the antibody
(\numrange{39}{71} clashing residues). Crucially this is \emph{not} a selection artifact: across all
\num{640} hotspot designs the clash floor is \num{14} and not one falls below \num{10}, so an accessible
minimal-hotspot graft was never generated --- the scorer is choosing the best of a pool in which every
member clashes. The three discontiguous residues appear to force the scaffold into the antibody's own
footprint, whereas the contiguous epitope leaves room to build clear of it. So the hotspot arm stands as
a demonstrated limitation of the minimal-graft idea: presenting isolated conserved residues in native
geometry seems to demand occupying the surface the antibody needs. We ship its top \num{10} anyway, as
that negative result. This arm is a self-contained component (\texttt{dp4\_8vdl/}); its \num{20} top
designs are merged into the library alongside the others at assembly.

\subsection{Assembly into one synthesis file}\label{sec:assembly}
All seven components (C1--C6 plus the 8VDL arm) share the constant 103-mer construct and the
eight-column annotated format (Section~\ref{sec:tiling}), so the final step concatenates them into a
single ordered library file with global \texttt{library\_member} numbering --- a named, documented
pipeline stage (\texttt{06\_library}, \texttt{scripts/stage06\_assemble.py}) so the shipped file is
regenerable from its components. This is done: at the shipped depth (top-\num{20} for C1/C2,
top-\num{10} for C3) assembly emits \num{14203} episcaf constructs, each with a unique
\texttt{library\_member} and \texttt{design\_ID}, all \num{103} residues.

\paragraph{The minibinder arm, and one assay for two projects.} Folded in at the same step are
\num{21759} \emph{minibinders}: de-novo designed binders against the same PfEMP1/EPCR target as the 8VDL
scaffolds, generated by a separate effort (LatentX) and filtered by their own criteria --- the passing
subset of $\sim$\num{484000} generations. They are not episcaf designs and were never scored on our
epitope-fidelity or accessibility axes, so their episcaf metric columns are blank; conversely they carry
their \num{13} native \texttt{lx\_*} metrics (per-residue and interface pLDDT, PAE, RMSD, \texttt{iptm},
target hotspots) on their own rows, preserved for later analysis. They are in the same file, and the same
synthesis order, for a practical reason: PepSeq builds one pooled oligo library per assay, so combining
two projects that target the same antigen into a single order shares the synthesis and the downstream
readout rather than paying for two. The combined file is \num{35962} rows (\num{14203} episcaf
$+$ \num{21759} minibinder); the \num{14203} is what the scoring and selection above concern, while the
encoding and order file cover all \num{35962}.

\paragraph{The cull to \num{35962}.} The library was first assembled at \num{37083} and then, after
review, cut to the shipped \num{35962} by four changes, none of which required re-encoding (every
surviving peptide keeps the oligo already computed for its sequence, matched back by sequence). (i) The
C6 controls were scaled from a top-\num{20} to a top-\num{15} base per antibody --- C6 was the single
largest block for what are controls --- keeping all three control flavors but recovering $\sim\!\num{775}$
peptides. To preserve the exact encoded sequences we \emph{filtered} the built C6 rather than rebuilding
it (rebuilding re-randomizes the scaffold-disruption controls, whose hexamer placement is drawn from a
global RNG). (ii) \num{286} peptides picked by more than one component were de-duplicated --- chiefly a C1
top design also drawn by C5's metric-space sample of the same pool --- so each peptide is ordered once.
(iii) The \num{60} designs from three single-island targets with a \num{2}-residue island were dropped:
two residues give too few epitope C$\alpha$ pairs to define the design-to-native fit, so their
accessibility (both the real clash and the cylinder) is uncomputable, and they had shipped with no
accessibility signal. (iv) A cosmetic naming artifact in \texttt{design\_ID} (a doubled RFdiffusion
backbone name from the AlphaFold3 output-directory convention) was collapsed. The final composition is
in Table~\ref{tab:dp4components}.

Two conventions hold across every row. First, \texttt{designedSequence} is always the full 103-mer in
\emph{EPITOPEscaffold} casing --- epitope residues uppercase, scaffold lowercase --- including the
linear controls, where the 30-mer tile is uppercase and the \texttt{GSGA} filler and \texttt{ENLYFQGA}
TEV site are lowercase; \texttt{sequence} carries the plain uppercase 103-mer that is synthesized.
Second, the eight annotation columns are followed by the full metric and scoring set, carried rather than
condensed: eight metric columns (epitope and overall RMSD; the PAE decomposition --- epitope, scaffold,
and global mean PAE; \texttt{ptm}; and both clash flavors, the real AF3 clash and the cylinder surrogate)
and four scoring columns (\texttt{composite}, \texttt{rank\_in\_group}, \texttt{is\_global\_pass}, and the
C2 \texttt{island\_index}), with a further \num{13} \texttt{lx\_*} columns appended on the minibinder rows
--- \num{33} columns in all. Every value is left blank where a design has no such measurement rather than
imputed: C3 has no AF3 clash (no antibody), 8VDL no cylinder (it has the real Fab), C5 no composite (it is
a metric-space sample, not composite-ranked), and C4 and C6 none of the fold metrics, since the linear
tiles and the mutants were never folded. The remaining mechanics --- recovering each design's sequence
--- follow.

\paragraph{Recovering the design sequences (case encoding).} Assembly and the C6 controls both need,
for each selected design, its full amino-acid sequence with the epitope residues marked --- stored as
a \emph{case-encoded} string (epitope uppercase, scaffold lowercase, so each contiguous uppercase run
is one island). This has to be rebuilt: our selection tables identify designs by a hash token, but the
epitope-position annotation lives in the design table \texttt{dp2} under a different key that was not
carried forward, so the token$\to$design mapping is missing for all but the assayed subset. We recover
it through the one key that always survives --- the design's own sequence: for each selected design we
read its chain-A sequence from its structure, match it to its \texttt{dp2} row, and uppercase the
epitope chunk-span positions (\texttt{scripts/case\_encode\_selected.py}, run once over the selected
set on the cluster). This yields the case-encoded sequences both C6 and assembly consume, and writes
the mapping back out so it need not be rebuilt again.

\subsection{From peptides to a DNA order: oligo encoding}\label{sec:oligo}

Everything up to this point is a library of \emph{peptides}. The assay does not receive peptides,
though --- it receives DNA. PepSeq builds each displayed peptide by transcribing and translating a DNA
oligo, so the final step converts all \num{35962} 103-mers of the combined library (the \num{14203}
episcaf constructs and the \num{21759} minibinders, encoded together into one assay) into the DNA that
codes for them and writes
the file a synthesis vendor (Twist) can manufacture. This subsection describes what that step does and
why, since it is the one part of the pipeline that crosses from computation into the wet lab; the
operational recipe and its pitfalls live in \texttt{episcaf\_pipeline/oligo\_encoding/README.md}.

\paragraph{Why encoding is a choice, not a lookup.} Translating a peptide to DNA is not a table
lookup, because the genetic code is degenerate: most amino acids are spelled by several
interchangeable codons, so a 103-residue peptide has an astronomical number of DNA sequences that all
translate to it. Those sequences are \emph{not} equally good to synthesize --- they differ in GC
content, in secondary structure, and in how faithfully the vendor can build them --- so ``encode the
library'' means \emph{choose} one good DNA spelling per peptide. We use the LadnerLab
\texttt{oligo\_encoding} tools (\url{https://github.com/LadnerLab/Library-Design}) with the same
updated codon-weight table as DP3
(\texttt{/tgen\_labs/Immunology/ekelley/DP3\_PepSeq\_Library\_Design/new\_codon\_weights/}), in two
stages: a sampler proposes many candidate encodings per peptide, weighted by the codon table and
targeting a GC ratio near \num{0.55}; then a pretrained model scores the candidates and keeps the
single best one. The result is one 309-nucleotide coding sequence per peptide ($103 \times 3$).

\paragraph{Adapters: the constant handles that make a pooled library usable.} The library is
synthesized as one pool --- tens of thousands of distinct oligos mixed in a single tube, each present
in trace amount. To do anything with that pool (amplify it, add the sequencing and display machinery,
read it back out) one needs a fixed sequence to grab onto, but every member differs in the middle.
The solution is to flank \emph{every} oligo with the same short constant sequences, one on each end,
called \emph{adapters} or handles. Because they are identical across all members, a single universal
primer pair binds them and amplifies the whole pool at once regardless of the variable coding sequence
between; the same handles are where the assay's display and readout elements attach. An adapter is
therefore a primer \emph{binding site}, not a primer itself --- the primer is a separate reagent, the
adapter's reverse complement.

Each final oligo is \texttt{[5\({}^\prime\) adapter]} $+$ \texttt{[309-nt coding sequence]} $+$
\texttt{[3\({}^\prime\) adapter]}. The adapters are fixed by the assay's primers rather than by our
design, so we pin them explicitly --- the \num{20}-mers DP3 shipped, giving a \num{349}-nt oligo ---
rather than leave them to a tool default.

That pinning is not fussiness. The flag that sets the adapters is a \emph{defaulted} one, and the two
installations of the encoder in circulation ship different defaults --- \num{19}-mers upstream against
the \num{20}-mers DP3 actually synthesized --- so building from a fresh clone silently yields oligos two
bases short, with nothing in the log to say so. The \num{19}-mer is the full primer footprint and sits
inside the \num{20}-mer, so the shorter form would very likely still prime; the objection is not that
\num{347} is known to break but that it is an unexplained deviation from the only construct known to
work, on an order that cannot be taken back. A defaulted flag in someone else's tool is a dependency on
which copy of the tool you happen to be running, and that is worth pinning wherever it reaches a physical
artifact.

\paragraph{The order file.} The output is the file sent to the vendor: two columns, a name and the
adapter-flanked DNA to synthesize, one row per library member. Because a synthesis order is expensive
and irreversible, we emit it through a checker (\texttt{scripts/stage07\_order\_file.py}) that verifies
every row before it goes out. The central check is a \emph{round-trip}: the encoder \emph{back-translates}
each peptide to DNA (choosing codons), so we \emph{translate} that DNA back through the standard genetic
code and require the result to equal the original peptide exactly, for all \num{35962} members. This
catches any codon-table, framing, or bookkeeping error at the one place it would otherwise pass silently
--- an encoding that is malformed but plausible-looking still fails the moment it does not read back to
its own peptide. The checker also confirms the adapters and the expected oligo length on every row. The
encoder itself is external tooling run on the cluster, so the computational pipeline ends at the
annotated peptide file and this checked order file is the deliverable that hands off to synthesis.

The full library has now been through it. All \num{35962} peptides encoded, every oligo \num{349}
nucleotides with the \num{20}-mer adapters on both ends, every core translating back to exactly the
peptide it claims to encode, one encoding per peptide and none missing
(\texttt{data/libraries/dp4\_order\_file.csv}).
